\section{Bilinear Embeddings}\label{a:ring_switch_remco}

Given a field $\mathbb F$ and a degree-$d$ extension $\mathbb E$. Our goal is to commit to a vector $\mathbf{w} \in \mathbb{F}^{n d}$ so that we can later use it to prove an inner product claim of the form $\langle \mathbf{w}, \mathbf{a}\rangle_\mathbb{F} = h$ for some $\mathbf{a} \in \mathbb{F}^{n d}$ and $h \in \mathbb{F}$. However, we only have an efficient inner product commitment scheme to commit to vectors $\widetilde{\mathbf{w}} \in \mathbb{E}^{n}$ and open as $\langle \widetilde{\mathbf{w}}, \widetilde{\mathbf{a}}\rangle_\mathbb{E} = \widetilde{h}$ for some $\widetilde{\mathbf{a}} \in \mathbb{E}^{n}$ and $\widetilde{h} \in \mathbb{E}$.

To do this, we seek a triple of maps $(W, A, H)$ that satisfy:
\[ 
\langle \mathbf{w},\mathbf{a}\rangle_{\mathbb{F}}
=
H\left(\langle W(\mathbf{w}), A(\mathbf{a}) \rangle_{\mathbb{E}} \right).
\]

\subsection{Background}

\subsubsection{Perfect pairings}
In the following let $\mathbb{E}/\mathbb{F}$ be a finite field extension of degree $d$. If we disregard multiplication in $\mathbb{E}$, then $\mathbb{E}$ is a $d$-dimensional vector space over $\mathbb{F}$, hence isomorphic to $\mathbb{F}^d$. Let $\mathbf{e}_1,\ldots,\mathbf{e}_d$ denote the standard basis of $\mathbb{F}^d$.
Recall that a function $\mu:\mathbb{E}\times \mathbb{E} \rightarrow \mathbb{F}$ is a \emph{bilinear form}  if it is linear in each argument separately. It is \emph{symmetric} if $\mu(x, y) = \mu(y, x)$ for all $x,y \in \mathbb{E}$.
For each $x \in \mathbb{E}$, the partial application in the first argument defines a linear form
\[
\mu_{x} : \mathbb{E} \rightarrow \mathbb{F} \quad y \mapsto \mu(x, y).
\]
Thus $\mu_{x} \in \mathbb{E}^*$ where $\mathbb{E}^* = \operatorname{Hom}_\mathbb{F}(\mathbb{E}, \mathbb{F})$ denotes the dual vector space. The bilinear form therefore induces an $\mathbb{F}$-linear map
\[
\Phi_\mu:\mathbb{E}\to\mathbb{E}^* \qquad x\mapsto\mu_x.
\]
We call $\mu$ \emph{nondegenerate}, or equivalently a \emph{perfect pairing}, if $\Phi_\mu$ is an isomorphism. Or equivalently for finite spaces, \albert{add a comment that injectivity implies surjectivity when the vector spaces have finite dimension?} % Remco: Like this?
\[
\mu(x,y)=0\text{ for all }y\in\mathbb E \quad\Longrightarrow\quad x=0.
\]
Since $\mu$ is a perfect pairing, every basis $(\beta_1,\ldots,\beta_d)$ of $\mathbb{E}$ has a unique dual basis $(\beta_1^\vee,\ldots,\beta_d^\vee)$ satisfying $\mu(\beta_i, \beta_j^\vee) = \delta_{ij}$.

\john{It would be good to explain how you would typically derive $\beta_j^\vee$}

\subsubsection{Succinct vectors}

Given a field $\mathbb{F}$, a vector $\mathbf{a} \in \mathbb{F}^n$ for $n = 2^k$ is called \emph{succinct} if it admits a tensor decomposition
\begin{equation}
\begin{bmatrix}
m_{00} \\ m_{01}
\end{bmatrix}
\otimes
\begin{bmatrix}
m_{10} \\ m_{11}
\end{bmatrix}
\otimes \dots \otimes
\begin{bmatrix}
m_{(k-1)0} \\ m_{(k-1)1}
\end{bmatrix}
\end{equation}
where typically $m_{ij} \in \mathbb{F}$. The multilinear extension of $\mathbf{a}$, viewed as a function on $\{0,1\}^k$, then becomes
\begin{equation}
\widetilde{a}(\mathbf{r}) = \prod_{i=0}^{k-1} \bigl( (1 - r_i) m_{i0} + r_i m_{i1} \bigr) \text{.}
\end{equation}
which can be evaluated in $O(k)$ field operations.

More generally, the coefficients $m_{ij}$ need not themselves be field elements. It suffices that they belong to an associative $\mathbb{F}$-algebra in which every complete product $m_{0b_0}m_{1b_1}\cdots m_{(k-1)b_{k-1}}$ evaluates to an element of $\mathbb{F}$. Consequently, every vector entry and every value of the multilinear extension remains an element of $\mathbb{F}$. As the $m_{ij}$ no longer need to be commutative we must therefore also fix a permutation of the variables $\mathbf{r}$.

Since every finite-dimensional associative algebra admits a matrix representation, we may, without loss of generality, assume that the objects $m_{ij}$ are matrices. In fact, it suffices to take $m_{ij}\in\mathbb{F}^{w_i\times w_{i+1}}$ with boundary dimensions $w_0=w_k=1$. We call $\max w_i$ the \emph{width} of the representation, which is known as a \emph{matrix branching program} and is closely related to matrix product states and weighted finite automatons.

The canonical example of a succinct vector is the multilinear evaluation vector $a_i = \mathrm{eq}(\mathbf{r}, \operatorname{bits}(i))$ obtained by taking $m_{i0} = 1 - r_i$ and $m_{i1} = r_i$. Another classical example is the geometric progression $(1, x, x^2, \dots)$ obtained by taking $m_{i0} = 1$ and $m_{i1} = x^{2^i}$. An illustrative non-scalar example is the counting vector $(0, 1, 2, \dots)$ represented by
\[
\begin{aligned}
m_{0b} &= \begin{bmatrix}
2 & b    
\end{bmatrix} &
m_{ib} &= \begin{bmatrix}
2 & b \\
0 &  1
\end{bmatrix} &
m_{(k-1)b} &= \begin{bmatrix}
b \\
1
\end{bmatrix}\text{.}
\end{aligned}
\]




\subsection{Structure}
We are now equipped to prove a theorem on the structure of embeddings. Under the assumption that the output map is $\mathbb{F}$-linear, we show that for $n=1$ the space of possible embeddings is very structured.

\begin{theorem}[Structure of bilinear embeddings]
Let $W,A:\mathbb{F}^d\rightarrow\mathbb{E}, H:\mathbb{E}\rightarrow\mathbb{F}$ be maps such that for all
$\mathbf{w},\mathbf{a}\in\mathbb{F}^d$,
\[
\langle \mathbf{w},\mathbf{a}\rangle_{\mathbb{F}}
=
H\bigl(W(\mathbf{w})A(\mathbf{a})\bigr).
\]
Suppose that $H$ is $\mathbb{F}$-linear. Then:
\begin{enumerate}
    \item $W$ and $A$ are $\mathbb{F}$-linear isomorphisms,
    \item $\mu_H(x,y) = H(xy)$ is a perfect pairing,
    \item $\left(W(\mathbf{e}_1),\ldots,W(\mathbf{e}_d)\right)$ and $\left(A(\mathbf{e}_1),\ldots,A(\mathbf{e}_d)\right)$ are dual bases with respect to $\mu_H$.
\end{enumerate}

Conversely, if $W$ and $A$, and $H$ are $\mathbb{F}$-linear, $\mu_H$ is a perfect pairing and the images of the standard basis under $W$ and $A$ are dual bases with respect to $\mu_H$, then the identity
\[
\langle \mathbf{w},\mathbf{a}\rangle_{\mathbb{F}}
=
H\bigl(W(\mathbf{w})A(\mathbf{a})\bigr).
\]
holds for all $\mathbf{w},\mathbf{a}\in\mathbb{F}^d$.
\end{theorem}

\begin{proof}[Proof of the forward direction]
First, observe that $
\mu_H(W(\mathbf e_i),A(\mathbf e_j))
= \langle \mathbf e_i,\mathbf e_j\rangle_{\mathbb F}
= \delta_{ij}
$. Taking $i=j$ shows $H$ is nonzero. Since $H$ is $\mathbb{F}$-linear, $\mu_H$ is a symmetric bilinear form, we show that it is nondegenerate. Let $x\in\mathbb E$ be nonzero. Since multiplication by $x$ is a bijection of $\mathbb E$ and $H$ is nonzero, there exists $y\in\mathbb E$ such that $\mu_H(x,y)=H(xy)\neq 0$. Thus $\mu_H$ is nondegenerate, and hence a perfect pairing.

Next, we show that $W$ and $A$ are bijective. Suppose
$W(\mathbf w_1)=W(\mathbf w_2)$. Then for every
$\mathbf a\in\mathbb F^d$,
\[
\langle \mathbf w_1-\mathbf w_2,\mathbf a\rangle_{\mathbb F}
= \mu_H(W(\mathbf w_1) - W(\mathbf w_2), A(\mathbf a))
=0.
\]
By nondegeneracy of the standard inner product on $\mathbb F^d$, this
implies $\mathbf w_1=\mathbf w_2$. Hence $W$ is injective. Since
$\mathbb F^d$ and $\mathbb E$ have the same finite cardinality, $W$ is
bijective. The same argument shows that $A$ is bijective.

We now prove that $W$ is $\mathbb F$-linear. For
$\mathbf w_1,\mathbf w_2,\mathbf a\in\mathbb F^d$, linearity of $H$
gives
\[
\begin{aligned}
&\mu_H\bigl(
W(\mathbf w_1+\mathbf w_2)-W(\mathbf w_1)-W(\mathbf w_2),
A(\mathbf a)
\bigr)\\
&=
\langle \mathbf w_1+\mathbf w_2,\mathbf a\rangle_{\mathbb F}
-
\langle \mathbf w_1,\mathbf a\rangle_{\mathbb F}
-
\langle \mathbf w_2,\mathbf a\rangle_{\mathbb F}\\
&=0.
\end{aligned}
\]
Sincle $A$ is surjective, the second argument ranges over all of
$\mathbb E$. Nondegeneracy of $\mu_H$ therefore implies
$W(\mathbf w_1+\mathbf w_2) = W(\mathbf w_1)+W(\mathbf w_2)$. Similarly, for $\lambda\in\mathbb F$,
\[
\begin{aligned}
&\mu_H\bigl(
W(\lambda\mathbf w)-\lambda W(\mathbf w),
A(\mathbf a)
\bigr)\\
&=
\langle \lambda\mathbf w,\mathbf a\rangle_{\mathbb F}
-
\lambda\langle \mathbf w,\mathbf a\rangle_{\mathbb F}\\
&=0.
\end{aligned}
\]
Again using surjectivity of $A$ and nondegeneracy of $\mu_H$, we obtain $W(\lambda\mathbf w)=\lambda W(\mathbf w)$. Hence $W$ is $\mathbb F$-linear. By symmetry, $A$ is also
$\mathbb F$-linear. Since they are also bijective, $W$ and $A$ are linear isomorphisms. From this it follows that $\left(W(\mathbf e_1),\ldots,W(\mathbf e_d)\right)$ and $
\left(A(\mathbf e_1),\ldots,A(\mathbf e_d)\right)$ are bases of $\mathbb E$. Finally, from $\mu_H(W(\mathbf e_i),A(\mathbf e_j)) = \delta_{ij}$ it follows they are dual with respect to $\mu_H$.
\end{proof}

\begin{proof}[Proof of the converse direction.] Can be directly shown using a basis expansion. 
Write $\mathbf w=\sum_{i=1}^d w_i\mathbf e_i$ and $\mathbf a=\sum_{j=1}^d a_j\mathbf e_j$
Since $W$ and $A$ are $\mathbb F$-linear, $W(\mathbf w) = \sum_{i=1}^d w_iW(\mathbf e_i)$ and $A(\mathbf a) = \sum_{j=1}^d a_jA(\mathbf e_j)$. Therefore,
\[
\begin{aligned}
H\bigl(W(\mathbf w)A(\mathbf a)\bigr)
&=
\mu_H\bigl(W(\mathbf w),A(\mathbf a)\bigr)\\
&=
\sum_{i,j}
w_i a_j
\mu_H\bigl(W(\mathbf e_i),A(\mathbf e_j)\bigr)\\
&=
\sum_{i,j}
w_i a_j\delta_{ij}\\
&=
\sum_iw_ia_i\\
&=
\langle\mathbf w,\mathbf a\rangle_{\mathbb F}.
\end{aligned}
\]
\end{proof}

The $n=1$ case extends naturally to $n>1$ by doing blockwise packing, i.e. we embed using $W' = \mathrm{I}_n \otimes W$ and $A' = \mathrm{I}_n \otimes A$ and keep the same $H$. Let $\mathbf{w} = (\mathbf{w}_1, \dots, \mathbf{w}_n)$ and $\mathbf{a} = (\mathbf{a}_1, \dots, \mathbf{a}_n)$ then
\[
\begin{aligned}
H\bigl( \langle W'(\mathbf{w}), A'(\mathbf{a}) \rangle_{\mathbb{E}} \bigr)
&= H\left( \sum_{i=1}^n W(\mathbf{w}_i) A(\mathbf{a}_i) \right) \\
&= \sum_{i=1}^n H\bigl( W(\mathbf{w}_i) A(\mathbf{a}_i) \bigr) \\
&= \sum_{i=1}^n \langle \mathbf{w}_i,\mathbf{a}_i \rangle_{\mathbb{F}} \\
&= \langle \mathbf{w},\mathbf{a} \rangle_{\mathbb{F}}.
\end{aligned}
\]
Going forward, we use $W$ and $A$ also for their blockwise extensions when the meaning is clear. We do not consider embeddings that mix coordinates across different blocks of size $d$, as such embeddings are unlikely to be useful in practice.

We have shown that every $n=1$ embedding with a linear output map is determined by a choice of a nonzero linear map $H$ and linear isomorphism $W$. The map $W$ then determines the first basis $\beta_i=W(\mathbf e_i)$ and the map $A$ is then forced to map to the $\mu_H$-dual basis.

The linearity assumption on the output map is necessary. Indeed, composing a linear embedding with a multiplicative bijection $g:\mathbb E\to\mathbb E$ generally produces an embedding that is nonlinear. Whether every nonlinear embedding arises from such a multiplicative reparameterization is a separate question, which we do not need here.

An important property of bilinear embeddings is that they preserve succinctness if the block size aligns with the tensor decomposition:

\begin{theorem}[Preservation of succinctness]\label{suc-embed}
Given $d = 2^k$, $n = 2^l$ and a $\mathbb{F}$-linear blockwise map $A = \mathrm{I}_n \otimes A_0$ for some isomorphism $A_0 : \mathbb{F}^d \to \mathbb{E}$, then if $\mathbf{a}$ \albert{$\mathbf{a} \in \FF^{d n}$} is succinct then so is $A(\mathbf{a})$. In particular if
\[
\mathbf{a} = \bigotimes_{i=0}^{k+l-1} 
\begin{bmatrix}
m_{i0} \\ m_{i1}
\end{bmatrix}
\]
where the $m_{ij}$ can be scalars or $w_i \times w_{i+1}$ matrices, then
\[
A(\mathbf{a}) = 
\left(
\bigotimes_{i=0}^{l-1} 
\begin{bmatrix}
m_{i0} \\ m_{i1}
\end{bmatrix}
\right) \cdot
(\mathrm{I}_{w_l} \otimes A_0)\left(
\bigotimes_{i=l}^{l+k-1} 
\begin{bmatrix}
m_{i0} \\ m_{i1}
\end{bmatrix}
\right)\text{.}
\]
where the second factor is column vector in $\mathbb{E}^{w_l}$ which can be absorbed into the matrices $m_{(l-1)b}$, bringing the representation back into canonical form.
\end{theorem}

\subsection{Constructions}

The set of nonzero linear output maps $H$ is $\mathbb{E}^* \setminus \{0\}$ which is isomorphic to $\mathbb{E}^\times$. The set of linear isomorphism $W$ is isomorphic to $\mathrm{GL}_d(\mathbb{F})$, hence the total number of bilinear embeddings is $|\mathbb{E}^\times| \cdot |\mathrm{GL}_d(\mathbb{F})|$. We'll now look at specific choices and their properties.

\subsubsection{Trace embedding}

Let $q = |\mathbb F|$, the $q$-power Frobenius map is $\varphi:\mathbb E\to\mathbb E$ given by $x\mapsto x^q$. It is an $\mathbb F$-linear field automorphism of order $d$, and its
fixed field is precisely $\mathbb F$:
\[
\mathbb E^\varphi
= \{x\in\mathbb E:\varphi(x)=x\}
= \mathbb F.
\]
The field trace is the $\mathbb F$-linear map $\operatorname{Tr}_{\mathbb E/\mathbb F}:\mathbb E\to\mathbb F$ defined by
\[
\operatorname{Tr}_{\mathbb E/\mathbb F}(x)
= \sum_{i=0}^{d-1}\varphi^i(x)
= \sum_{i=0}^{d-1}x^{q^i}.
\]
To see that this sum is indeed in $\mathbb{F} \subseteq \mathbb{E}$ observe the sum is fixed by $\varphi$, since
\[
\varphi\left( \sum_{i=0}^{d-1}\varphi^i(x) \right)
= \sum_{i=1}^{d}\varphi^i(x)
= \sum_{i=0}^{d-1}\varphi^i(x),
\]
where $\varphi^d(x)=\varphi^0(x)$. Hence the trace lies in the fixed field $\mathbb F$.

This makes $\operatorname{Tr}_{\mathbb E/\mathbb F}$ a natural choice of output map $H$. More generally, every $\mathbb F$-linear map $H:\mathbb E\to\mathbb F$ is uniquely of the form $H(x)=\operatorname{Tr}_{\mathbb E/\mathbb F}(\alpha x)$ for some $\alpha\in\mathbb E$. Consequently, the nonzero output maps are parameterized by $\mathbb E^\times$.

The Galois orbit embedding chooses the basis $\beta_i=\varphi^i(\gamma)$ which has dual $\beta_i^\vee=\varphi^i(\gamma^\vee)$.

\ulrich{%
This sounds to me very much related to Galois Fribinius on an inner product with a general covector. 
You would end up with the multi-claim of the inner products of the packed poly with the Galois orbit of the covector.
This is how I view it:
Given $\vec p_1, \ldots, \vec p_d \in \mathbb F^n$, and $\vec q_1, \ldots, \vec q_d \in \mathbb F^n$, we are interested in 
\[
\sum_{k=1}^d \langle \vec p_k, \vec q_k\rangle,
\]
right?

The packing map sends the coeff vectors $\vec p_1,\ldots, \vec p_d$ to $\vec p = \beta_1\cdot \vec p_1 + \ldots + \beta_d\cdot \vec p_d$, where $\beta_1,\ldots,\beta_d$ is an $\mathbb F$-basis of $\mathbb E$. 
Choose the element $a_k$ in $\mathbb E$ which extracts the $k$-th coordinate
\[
\vec p_k = \Phi_k(\vec p) = \operatorname{Tr}_{\mathbb E/\mathbb F}(a_k\cdot \vec p).
\]
Then 
\begin{align*}
\langle \vec p_k, \vec q_k\rangle 
&= \langle \sum_{i=0}^d \varphi^i(a_k\cdot \vec p), \vec q_k\rangle 
\\
&= \sum_{i=0}^d  \varphi^{i}(a_k) \cdot \varphi^i(\langle \vec p,  \vec q_k\rangle),
\end{align*}
since $\vec q_k$ is over $\mathbb F$.
This reduces our inner product claim to that of the packed poly $\vec p$ with the unpacked polys $\vec q_1,\ldots, \vec q_k$.

But does that match our use case?
We are rather interested in $\langle \vec p_k, \vec v\rangle$ for some $\vec v\in \mathbb E^n$.
In this case, 
\begin{align*}
\langle \vec p_k, \vec v \rangle 
&= \langle \sum_{i=0}^d \varphi^i(a_k\cdot \vec p), \vec q_k\rangle 
\\
&= \sum_{i=0}^d  \varphi^{i}(a_k) \cdot \varphi^i(\langle \vec p,  \varphi^{-i}(\vec v)\rangle),
\end{align*}
and we end up with inner products of the packed poly with the Galois orbit of $\vec v$. 
}


\subsubsection{Coefficient projection}

In the coefficient projection we assume a concrete representation of $\mathbb{E}$ as $\mathbb{F}[X]/(f(X))$ for some monic degree $d$ polynomial $f$ irreducible over $\mathbb{F}$. We assume $f$ is chosen for implementation efficiency: it is sparse with mostly low order terms. Further we represent members of $\mathbb{E}$ in the monomial basis $(1, X, X^2, \dots, X^{d-1})$. We can then represent $W$ and $A$ by invertible matrices in $\mathbb{F}^{d \times d}$ and $H$ by a vector in $\mathbb{F}^d$. For implementation efficiency we seek an embedding with sparse $W$ and $A$. It turns out taking the most efficient choices for $H$ and $W$ also gives efficient $A$:

\begin{theorem}
Write  $f(X)= f_d X^d + f_{d-1} X^{d-1} + \hdots + f_1 X + f_0$ so that $f_d =1$ and $f_0 \neq 0$. Take $W=\mathrm{I}_d$ and $H=(1, 0, \dots, 0)$. Then the dual-basis embedding $A$ is given by
\[
A_{ij}=
\begin{cases}
1,&i=j=0,\\
-f_{i+j}/f_0,&i,j\ge1\text{ and }i+j\le d,\\
0,&\text{otherwise}.
\end{cases}
\]
\end{theorem}

In particular, $A$ is a bordered upper anti-triangular Hankel matrix. For $d=5$ it looks like
\[
A= \frac{1}{f_0}
\begin{pmatrix}
f_0 & 0 & 0 & 0 & 0\\
0 & -f_2 & -f_3 & -f_4 & -1\\
0 & -f_3 & -f_4 & -1 & 0\\
0 & -f_4 & -1 & 0 & 0\\
0 & -1 & 0 & 0 & 0
\end{pmatrix}.
\]
Notably $f_1$ does not appear and for $k>1$ the $f_k$ coefficient appears $k-1$ times. This perfectly matches the usual implementation preference for sparse reduction polynomials with mostly low-order terms.

When doing inner product openings, the $1/f_0$ factor can be moved out of the inner product. 

For the 127-bit field $\mathbb{F}_2[X]/(X^{127} + X + 1)$ the $A$ map is simply a reversal of the $i>0$ coordinates. For the GHASH field $\mathbb{F}_2[X]/(X^{128} + X^7 + X^2 + X + 1)$ the map additionally needs seven XOR corrections.

% Cite: https://shiftleft.com/mirrors/www.hpl.hp.com/techreports/98/HPL-98-135.pdf


\subsection{Extension openings}

Assume now that we wish to open a base-field witness $\mathbf{w} \in \mathbb{F}^{nd}$ against extension-field weights $\mathbf{a} \in \mathbb{E}^{n d}$, proving the claim $\langle \mathbf{w}, \mathbf{a}\rangle_\mathbb{E} = h$ for some $h \in \mathbb{E}$. We present a one-round batching protocol and a $log d$-round folding protocol.

\subsubsection{Batching protocol}
Given a bilinear embedding $(W, A, H)$ and a basis $(\beta_1, \dots, \beta_d)$ for $\mathbb{E}/\mathbb{F}$, we reduce the extension-field claim by decomposing into basis coordinates and batching     the resulting claims with a random linear combination.

\begin{construction}[Batching protocol for extension opening]
Given the claim $\langle \mathbf{w}, \mathbf a\rangle_\mathbb{E}=h$, define $\mathbf{a}_i \in \mathbb{F}^{nd}$ such that 
\[
\mathbf{a} = \sum_{i=1}^d \mathbf{a}_i \beta_i \text{.}
\]
The parties proceed as follows.
\begin{enumerate}
    \item The prover sends $h_i = \langle W(\mathbf{w}), A(\mathbf{a}_i) \rangle_{\mathbb{E}}$  for $i=1,\ldots,d$.
    \item The verifier checks $h = \sum_{i=1}^{d} H(h_i) \beta_i$.
    \item The verifier samples an $\epsilon$-zero-evader $\boldsymbol{\rho} \sample\mathbb E^d$.
    \item The parties reduce the claim to $\langle W(\mathbf{w}), \mathbf a'\rangle_\mathbb{E}=h'$ where 
    \[
    h' = \sum_{i=1}^{d} \rho_i h_i,
    \qquad
    \mathbf{a}' = \sum_{i=1}^{d} \rho_i  A(\mathbf{a}_i) \text{.}
    \]
\end{enumerate}
This interactive oracle reduction is always complete, has round-by-round soundness error $\epsilon$, and proof size $d \log |\mathbb{E}| = d^2 \log |\mathbb{F}|$.
\end{construction}

\begin{theorem}[Succinctness of coordinate decomposition]
Given a succinct vector $\mathbf{a} \in \mathbb{E}^{nd}$ for $d= 2^k$ and a basis $(\beta_i)$ of $\mathbb{E}/\mathbb{F}$ then the $\beta$-coordinate decomposed vectors are also succinct. In particular if
\[
\mathbf{a} = \bigotimes_{i=0}^{k+l-1} 
\begin{bmatrix}
m_{i0} \\ m_{i1}
\end{bmatrix}
\]
where the $m_{ij}$ can be scalars or $w_i \times w_{i+1}$ matrices, then
\[
\mathbf{a}_i = \mathbf{u}_i \cdot
\left(
\bigotimes_{j=0}^{k+l-1} 
\begin{bmatrix}
\Phi (m_{j0}) \\ \Phi(m_{j1})
\end{bmatrix}
\right)
\cdot \mathbf{v}
\]
where $\Phi : \mathbb{E} \to \mathbb{F}^{d \times d}$ is the map that represents an $\mathbb{E}$-element $m$ by the matrix representing multiplication by $m$ in basis $(\beta_i)$, $\mathbf{u}_i$ is the one-hot vector with $u_{ij} = \delta_{ij}$, and $\mathbf{v}$ is the base $(\beta_i)$ representation of $1 \in \mathbb{E}$. $\Phi$ is applied elementwise to $w_i \times w_{i+1}$ matrices over $\mathbb{E}$ to produce $dw_i \times dw_{i+1}$ block matrices over $\mathbb{F}$. Absorbing the outer vectors in the matrices brings the representation back into canonical form, with the width increased by at most a factor $d$.
\end{theorem}

For practical implementations it makes sense to evaluate the branching program over $\mathbb{E}$ directly and then extract the coordinates in the basis $(\beta_i)$. This is equivalent to omitting the output selector $\mathbf{t}_i$ to keep the output dimension $w_0 = d$.

If $A = \mathrm{I}_n \otimes A_0$ is blockwise we can also apply theorem \ref{suc-embed} and compose it with the combination weights $\boldsymbol{\rho} \in \mathbb{E}^d$ to get a succinct representation for $\mathbf{a}'$:
\begin{equation}
\mathbf{a}' = 
\boldsymbol{\rho} \cdot
\left(
\bigotimes_{i=0}^{l-1} 
\begin{bmatrix}
\Phi (m_{i0}) \\ \Phi (m_{i1})
\end{bmatrix}
\right) \cdot
(\mathrm{I}_{d w_l} \otimes A_0)\left(
\bigotimes_{i=l}^{l+k-1} 
\begin{bmatrix}
\Phi (m_{i0}) \\ \Phi(m_{i1})
\end{bmatrix}
\right) \cdot \mathbf{v}
\text{.}
\end{equation}
The map $\Phi$ is an algebra embedding, so replacing each scalar $m_{ij}\in\mathbb E$ by $\Phi(m_{ij})$ leaves the branching-program computation unchanged, except that the values are now represented in the basis $(\beta_i)$ representation of $\mathbb E$. The only operations outside the embedded $\mathbb{E}$ computation are the blockwise embedding $A_0$ and the random linear combination $\boldsymbol{\rho}$. Thus, apart from the applications of $A_0$ and $\boldsymbol{\rho}$, evaluating $\widetilde{\mathbf a}'$ is the same computation as evaluating $\widetilde{\mathbf a}$, performed in the regular representation of $\mathbb E$. \remco{Something is off with the complexity. Should really check with a Sage implementation. Probably needs a factor $d$.}

The batch protocol is quadratic in $d$ in terms of $\mathbb{F}$-element operations and proof size. This is acceptable when $d$ is small, but binary fields where $d = 128$ the communication would be $2,048$ bytes. Avoiding this quadratic dependence on $d$ appears to require further rounds of interaction and more sophisticated decomposition.

\subsubsection{Folding protocol}
In the folding protocol we interactively reduce the claim $\langle \mathbf{w}, \mathbf{a}\rangle_\mathbb{E} = h$ through a sequence of embeddings $W_i$ where at round $i$ the claim is
\[
\begin{array}{ccc}
\langle W_i(\mathbf{w}), A_i(\mathbf{a})\rangle_\mathbb{E} = h_i
& W_i: \mathbb{F}^{nd} \to \mathbb{E}^{nd_i}
& A_i: \mathbb{E}^{nd} \to \mathbb{E}^{nd_i}
\end{array}
\]
Here $d = d_0 \ge \dots \ge d_k = 1$ is a series of folding dimensions. We start with $W_0$ and $A_0$ the identity and $h_0 = h$ so that the initial claim is indeed $\langle \mathbf{w}, \mathbf{a} \rangle_\mathbb{E} = h$. We also stipulate that the embeddings are blockwise such that $W_i = \mathrm{I}_n \otimes W_i'$ and $A_i = \mathrm{I}_n \otimes A_i'$. For simplicity we we will only consider the $n=1$ case from here on as $n>1$ follows by blockwise repetition.

The final map is $W_k : \mathbb{F}^d \to \mathbb{E}$ which is an $F$-linear isomorphism. Let $(\beta_i)$ be the implied basis such that 
\[
W_k(\mathbf x) = \sum_{i=1}^d x_i \beta_i \text{.}
\]

Working backwards from the final embedding we create a sequence
\[
R_i \circ W_{i+1} = W_i
\]
where each map $R_i$ has only $\left\lceil \frac{d_i}{d_{i+1}}\right\rceil$ Frobenius components.

\pagebreak
Every $\mathbb F$-linear map $T:\mathbb E^m\to\mathbb E^n$ can be viewed as being built from $\mathbb E$-linear maps and $\mathbb F$-linear scalar operations on $\mathbb E$. To recursively fold such maps, the scalar operations must be compatible with arbitrary $\mathbb E$-linear coordinate transformations. Concretely, we require an intertwining
\[
\beta(R\mathbf x)=\alpha(R)\beta(\mathbf x),
\]
where $\alpha$ acts entrywise on R. This property singles out semilinear scalar operators. Over finite fields, these are precisely the Frobenius maps $\sigma^i$ up to scalar multiples, and the Frobenius powers therefore provide the canonical scalar decomposition.

% https://arxiv.org/pdf/1211.5475

\begin{theorem}[Frobenius decomposition]
Let $\mathbb E/\mathbb F$ be a finite extension of degree $d$, and let $\sigma: \mathbb{E} \to \mathbb{E}$ be the Frobenius map $x \mapsto x^{|\mathbb{F}|}$. Then every $\mathbb F$-linear map $T: \mathbb{E}^m \to \mathbb{E}^n$ admits a unique decomposition
\[
T(\mathbf{x}) = \sum_{i=0}^{d-1} M_i \cdot (\mathrm{I}_m \otimes \sigma^i)(\mathbf{x})
\text{,} \qquad M_i \in \mathbb{E}^{n \times m} \text{.}
\]
\end{theorem}
We define $\operatorname{supp}_\sigma(T) := \{i : M_i \neq 0\}$.

\paragraph{Single folding round.}
Assume the current claim is
\[
\langle W_i(\mathbf w),A_i(\mathbf a)\rangle_{\mathbb E}=h_i.
\]
Let
\[
R_i:\mathbb E^{nd_{i+1}}\to\mathbb E^{nd_i}
\]
be a fixed $\mathbb F$-linear map satisfying
\[
R_i\circ W_{i+1}=W_i.
\]

By the Frobenius decomposition there exist unique matrices
\[
M_{i,0},\ldots,M_{i,d-1}
\in
\mathbb E^{nd_i\times nd_{i+1}}
\]
such that
\[
R_i
=
\sum_{j=0}^{d-1}
M_{i,j}\circ
(\mathrm I_{nd_{i+1}}\otimes\sigma^j).
\]

For every $j$, define the $\mathbb F$-linear map
\[
C_{i,j}
=
(\mathrm I_{nd_{i+1}}\otimes\sigma^{-j})
\circ
M_{i,j}^{\mathsf T}.
\]
Then, for every
$\mathbf x\in\mathbb E^{nd_{i+1}}$
and
$\mathbf b\in\mathbb E^{nd_i}$,
\[
\langle R_i(\mathbf x),\mathbf b\rangle_{\mathbb E}
=
\sum_{j=0}^{d-1}
\sigma^j\!\left(
\langle
\mathbf x,
C_{i,j}(\mathbf b)
\rangle_{\mathbb E}
\right).
\]

Applying this identity to
\[
\mathbf x=W_{i+1}(\mathbf w),
\qquad
\mathbf b=A_i(\mathbf a),
\]
gives
\[
\langle
W_i(\mathbf w),
A_i(\mathbf a)
\rangle_{\mathbb E}
=
\sum_{j=0}^{d-1}
\sigma^j\!\left(
\langle
W_{i+1}(\mathbf w),
C_{i,j}(A_i(\mathbf a))
\rangle_{\mathbb E}
\right).
\]

The prover sends the values
\[
g_{i,j}
=
\langle
W_{i+1}(\mathbf w),
C_{i,j}(A_i(\mathbf a))
\rangle_{\mathbb E},
\qquad
j=0,\ldots,d-1.
\]
The verifier checks
\[
h_i
=
\sum_{j=0}^{d-1}
\sigma^j(g_{i,j}),
\]
samples a random challenge
$\rho_i\in\mathbb E$,
and defines
\[
A_{i+1}
=
\left(
\sum_{j=0}^{d-1}
\rho_i^j C_{i,j}
\right)
\circ A_i.
\]
Finally, letting
\[
h_{i+1}
=
\sum_{j=0}^{d-1}
\rho_i^j g_{i,j},
\]
the next claim is
\[
\langle
W_{i+1}(\mathbf w),
A_{i+1}(\mathbf a)
\rangle_{\mathbb E}
=
h_{i+1}.
\]


We therefore seek an interactive reduction from a mixed claim to a packed claim
\[
\begin{array}{ccc}
\langle \mathbf{w}, \mathbf{a}\rangle_\mathbb{E} = h
&
\mathbf{w} \in \mathbb{F}^{nd}
&
\mathbf{a} \in \mathbb{E}^{nd}
\\[.8ex]
\Downarrow
\\[1.3ex]
\langle W(\mathbf{w}), \mathbf{a}'\rangle_\mathbb{E} = h'
&
W(\mathbf{w}) \in \mathbb{E}^{n}
&
\mathbf{a}' \in \mathbb{E}^{n}
\end{array}
\]
where $W: \mathbb{F}^{nd} \to \mathbb{E}^{n}$ is a fixed $\mathbb{F}$-linear embedding and $\mathbf{a}' \in \mathbb{E}^{n}$, $h' \in \mathbb{E}$ are determined by $\mathbf{a}$, $h$ and the transcript. Ideally $\mathbf{a} \mapsto \mathbf{a}'$ should also preserve useful properties like efficient computation of its multilinear extension.

Above we decomposed $\mathbf{a}$ with respect to a basis of $\mathbb{E}/\mathbb{F}$. Here we instead decompose the induced $\mathbb F$-linear functional on packed witnesses. Define
\[
T_{\mathbf{a}} : \mathbb{E}^n \to \mathbb{E} \qquad \mathbf{x} \mapsto \langle W^{-1}(
\mathbf{x}), \mathbf{a} \rangle_{\mathbb{E}}\text{.}
\]
This is an $\mathbb{F}$-linear functional on $\mathbb E^n$, but is not generally $\mathbb{E}$-linear. We'll first prove a representation lemma for $\mathbb{F}$-linear $\mathbb{E}\to\mathbb{E}$ maps and then a decomposition theorem.

\begin{theorem}
Let $\mathbb{F} \subseteq \mathbb{E}$ be finite fields with $q = |\mathbb{F}|$ and $[\mathbb{E}:\mathbb{F}] = d$,then every $\mathbb{F}$-linear map $T: \mathbb{E} \to \mathbb{E}$ is uniquely represented by a $q$-linearized polynomial
\[
P(X)=\sum_{j=0}^{d - 1} c_j X^{q^j} \qquad c_j\in\mathbb E \text{.}
\]
\end{theorem}
\begin{proof}
By interpolation every map $T: \mathbb{E} \to \mathbb{E}$ can be represented uniquely by a polynomial $P \in \mathbb{E}[X]$ of degree less than $|\mathbb{E}| = q^d$
\[
P(X)=\sum_{j=0}^{q^d - 1} a_j X^j \qquad a_j\in\mathbb E \text{.}
\]
From $\mathbb{F}$-linearity follows $P(0) = 0$ and hence $a_0 = 0$, further we have $P(X + Y) = P(X) + P(Y)$ and hence
\[
\begin{aligned}
0 &= P(X + Y) - P(X) - P(Y) \\
&= \sum_{j=1}^{q^d - 1} a_j (X + Y)^j - \sum_{j=1}^{q^d - 1} a_j (X^j + Y^j) \\
&= \sum_{j=1}^{q^d - 1} a_j \sum_{k=0}^{j} \binom{j}{k}X^k Y^{j-k} - \sum_{j=1}^{q^d - 1} a_j (X^j + Y^j)\\
&= \sum_{j=1}^{q^d - 1} a_j \sum_{k=1}^{j-1} \binom{j}{k}X^k Y^{j-k}
\end{aligned}
\]
Write $q=p^r$ with $p$ the characteristic of $\mathbb{F}$. For every $j$ with $a_j \neq 0$ we must have $\binom{j}{k} =0 \mod p$ for all $0<k<j$, by Lucas’s theorem this holds iff $j = p^s$ for some $s$. Hence we can reindex $P$ keeping only the non-zero coefficients
\[
P(X)=\sum_{j=0}^{rd - 1} b_j X^{p^j} \qquad b_j\in\mathbb E \text{.}
\]
From $\mathbb{F}$-linearity also follows $P(\lambda x) = \lambda P(x)$ for all $x \in \mathbb{E}$ and $\lambda \in \mathbb{F}$. For every $b_j \neq 0$ this requires $\lambda^{p^j} = \lambda$ for every $\lambda \in \mathbb{F}$. This holds exactly when $j = r s$ for some $s$. Reindexing $P$ again we get our claimed result
\[
P(X)=\sum_{j=0}^{d - 1} c_j X^{q^j} \qquad c_j\in\mathbb E \text{.}
\]
\end{proof}



\begin{theorem}[Frobenius decomposition]
Let $\mathbb{F} \subseteq \mathbb{K} \subseteq \mathbb{E}$ be finite fields with $[\mathbb{K}:\mathbb{F}] = m$, and let $\sigma(x) = x^{|\mathbb{F}|}$ denote the $\mathbb{F}$-Frobenius automorphism. Then every $\mathbb{F}$-linear map $T: \mathbb{E} \to \mathbb{E}$ admits a unique decomposition
\[
T(x) = \sum_{i = 0}^{m - 1} \sigma^i \circ T_i
\]
where each $T_i$ is $\mathbb{K}$-linear.
\end{theorem}
\begin{proof}
Write $q=|\mathbb F|$ and $d=[\mathbb E:\mathbb F]$, by the preceding theorem $T$ has a unique representation
\[
T(X)=\sum_{j=0}^{d - 1} c_j X^{q^j} \qquad c_j\in\mathbb E \text{.}
\]
Since $\mathbb K\subseteq\mathbb E$, we have $m\mid d$. Write each
index uniquely as $j=ms+i$, with $0\leq i<m$ and $0\leq s<d/m$. Grouping the terms by their residue modulo $m$ gives
\[
T(X)
=
\sum_{i=0}^{m-1}
\sum_{s=0}^{d/m-1}
c_{ms+i}X^{q^{ms+i}}.
\]
For each residue class $i$, define
\[
T_i(X)
=
\sum_{s=0}^{d/m-1}
c_{ms+i}^{q^{-i}}X^{q^{ms}},
\]
where $z\mapsto z^{q^{-i}}$ denotes the inverse of the automorphism $z\mapsto z^{q^i}$ of $\mathbb E$. Each $T_i$ is now a $q^m$-linearized polynomial hence applying the preceding theorem to the extension $\mathbb{K} \subseteq \mathbb{E}$, each $T_i$ is $\mathbb{K}$-linear.  Then
\[
(\sigma^i\circ T_i)(X)
= T_i(X)^{q^i}
= \sum_{s=0}^{d/m-1} c_{ms+i}X^{q^{ms+i}} \text{.}
\]
Therefore $T$ represented as
\[
T=\sum_{i=0}^{m-1}\sigma^i\circ T_i \text{.}
\]
Uniqueness follows immediately from the uniqueness of the $q$-linear polynomial representation.
\end{proof}

The decomposition theorem extends coordinatewise to maps $\mathbb E^n\to\mathbb E$. In particular, the functional $T_{\mathbf a}(\mathbf{x}) = \langle W^{-1}(\mathbf{x}), \mathbf{a} \rangle_{\mathbb{E}}$ admits a unique decomposition
\begin{equation}
T_{\mathbf a}
=
\sum_{i=0}^{m-1}\sigma^i\circ T_i,
\end{equation}
where each $T_i:\mathbb E^n\to\mathbb E$ is $\mathbb K$-linear. Thus the initial claim may be written as
\[
h = T_{\mathbf a}(W(\mathbf w))
= \sum_{i=0}^{m-1} \sigma^i\bigl(T_i(W(\mathbf w))\bigr).
\]

\begin{construction}[Frobenius fold]
Given the claim $T_{\mathbf a}(W(\mathbf w))=h$, the parties proceed as follows.
\begin{enumerate}
    \item The prover sends $v_i=T_i(W(\mathbf w))$  for $i=1,\ldots,m-1$.
    \item The verifier derives $v_0 = h-\sum_{i=1}^{m-1}\sigma^i(v_i)$.
    \item The verifier samples independent challenges $\rho_1,\ldots,\rho_{m-1}\sample\mathbb E$.
    \item The parties reduce to the claim $T'(W(\mathbf w))=h'$, where
    \[
    T' = T_0+\sum_{i=1}^{m-1}\rho_iT_i
    \qquad\text{and}\qquad
    h' = v_0+\sum_{i=1}^{m-1}\rho_iv_i.
    \]
\end{enumerate}
\end{construction}

\begin{construction}[Folding protocol for extension opening]
Choose a tower of intermediate fields $\mathbb F=\mathbb K_0 \subseteq\mathbb K_1 \subseteq\cdots \subseteq\mathbb K_t=\mathbb E$, and write $m_i=[\mathbb K_{i+1}:\mathbb K_i]$. Initially, $T_{\mathbf a}$ is $\mathbb K_0$-linear. At round $i$, apply the Frobenius fold with $\mathbb F=\mathbb K_i$, and $\mathbb K=\mathbb K_{i+1}$.  This reduces the current $\mathbb K_i$-linear functional to a $\mathbb K_{i+1}$-linear functional.

After $t$ rounds, the residual functional $T':\mathbb E^n\to\mathbb E$ is $\mathbb E$-linear. Consequently, there exists $\mathbf a'\in\mathbb E^n$ such that $T'(\mathbf x)=\langle\mathbf x,\mathbf a'\rangle_{\mathbb E}$. The final claim is therefore $\langle W(\mathbf w),\mathbf a'\rangle_{\mathbb E}=h'$, as required.
\end{construction}

Each round requires sending $m_i - 1$ elements of $\mathbb{E}$. For a binary tower where $m_i=2$, $\mathbb{F} = \mathbb{F}_2$, and $\mathbb{E} = \mathbb{F}_{2^{128}}$ we have $7$ rounds of one $\mathbb{E}$ element message each, for a total of $112$ bytes. Compare to the non-interactive protocol this reduces the communication complexity from $O(d^2)$ to $O(d \log d)$ elements of $\mathbb{F}_2$.

\paragraph{Towerless folding.}
Say $\mathbb{F} \subseteq \mathbb{E}$ admits no tower of subfields, for example $\mathbb{E} = \mathbb{F}_{2^{127}}$. Take again $d = [\mathbb{E}: \mathbb{F}]$ and decompose $T_\mathbf{a}$ as
\[
T_{\mathbf a} = \sum_{i=0}^{d-1}\sigma^i\circ T_i,
\]
We can still recursively fold using the protocol above with $m=2$ by grouping even and odd exponents of $\sigma^i$ (zero-padding as necessary), as there is no actual dependence on the subfield.

This folding progressively halves $d$ (rounding up) in the $\mathbb{E}^{nd}$ inner product. This induces a sequence of perfect pairings 
\[
(\mathbb F^{nd},\mathbb E^{nd}) \rightsquigarrow (\mathbb E^{nd/m},\mathbb E^{nd/m}) \rightsquigarrow (\mathbb K_i^{nd/m},\mathbb E^{nd/m})
\] 

\[
\langle
W_i(\mathbf w),
A_i(\mathbf a)
\rangle
=
\langle
\mathbf w,
\mathbf a
\rangle.
\]

\paragraph{}
The final result is a transcript dependent map $A: \mathbb{E}^{dn} \to \mathbb{E}^n$. If $W$ is blockwise with $\mathrm{I}_n \otimes W_0$ for some $W_0: \mathbb{F}^d \to \mathbb{E}$ then so is $A$ with $A_0: \mathbb{E}^d \to \mathbb{E}$.



